\begin{resumo}
Este trabalho apresenta o desenvolvimento de um chatbot voltado ao apoio acadêmico de estudantes da Universidade de Brasília (UnB), com base em uma abordagem centrada no usuário, fundamentada nos princípios do ITIL v4 e do Design Thinking. A proposta surge da necessidade de centralizar e simplificar o acesso a informações institucionais, frequentemente dispersas e de difícil assimilação. A solução consiste em um bot funcional, integrado a plataformas amplamente utilizadas pelos estudantes, como WhatsApp, Telegram e uma aplicação web complementar. A ferramenta foi desenvolvida a partir da estruturação de documentos oficiais da UnB em uma base de dados, permitindo ao chatbot responder dúvidas frequentes sobre prazos, procedimentos administrativos, disciplinas, TCCs, estágios, entre outros temas. Para validar a eficácia da solução, foi realizada uma pesquisa com discentes do campus do Gama (FCTE), cujos resultados contribuíram diretamente para a definição dos requisitos do sistema. Além disso, o bot oferece um canal de encaminhamento para atendimento humano, quando necessário, conectando os estudantes a setores institucionais como a secretaria acadêmica. A partir dos requisitos levantados, foi desenvolvido um Mínimo Produto Viável (MVP) que demonstra como a proposta pode promover maior autonomia e agilidade no acesso às informações, além de aprimorar a comunicação entre estudantes e a universidade.

 \vspace{\onelineskip}
    
 \noindent
 \textbf{Palavras-chave}:  chatbot. apoio acadêmico. ITIL v4. Design Thinking. UnB.
\end{resumo}
